\begin{abstract}
As a field, Evolutionary Computation (EC) is deeply entwined with its heuristic metaphor to biological evolution.
This framing has driven numerous bio-inspired innovations and embellishments --- some productive, and many less so.
Here, we argue that deeper value from bio-inspiration lies not in biological systems themselves, but in the fields of science they begat.
Because EC and natural evolution both instantiate heritable variation under selection, large bodies of theory and analytical methods from evolutionary biology and allied fields can be recruited to EC.
We review existing steps in this direction --- by application of population genetics, ecology, biodiversity science, adaptive landscapes, and evolutionary graph theory into EC.
We believe much untapped potential remains, however, and suggest promising opportunities to recruit inferential observation, replay experiments, coevolution, adaptive momentum, and phylogenetic network analysis.
% Notably, exchange between EC and bioscience has long been a two-way street; we also highlight instances of and opportunities for reciprocal exchange.
Such interdisciplinary awareness and fluency will be key to advancing evolutionary computation towards its full potential, and broadening its impact by also contributing knowledge back to bioscience.
\end{abstract}
